% ============================================
% LUC-Escolha
% Gerado de molambudos.md — build_miolo.py
% ============================================
\section*{\textbf{LUC-Escolha} --- A Escolha}
\markboth{LUC-Escolha --- A Escolha}{}

\noindent \textit{O que fazer com o diabo que você carrega}

\noindent \textbf{Abril de 2026 --- Barbacena}

\vspace{0.5em}\noindent\rule{\textwidth}{0.5pt}\vspace{0.5em}

\noindent Lúcia passou a noite no único hotel de Barbacena. Um quarto pequeno, mofo nas paredes, ventilador de teto que girava como um carrossel lento.

\noindent Ela colocou os dois diários sobre a mesa: o original (capa marrom, 62 anos de escrita) e o segundo (capa preta, escrito depois da morte).

\noindent Ela não abriu o segundo. Ainda não.

\noindent Em vez disso, sentou na cama e pensou.

\noindent Ela tinha três opções.

\noindent \textbf{Opção 1: Devolver o diário ao arquivo.}

\noindent Fingir que nada aconteceu. Deixar que outro pesquisador o encontrasse. O ciclo continuaria --- mas ela não seria responsável. Ou seria? Ela já tinha lido. A contaminação já tinha começado. Devolver o diário não apagaria os sintomas.

\noindent \textbf{Opção 2: Destruir o diário.}

\noindent Queimar. Rasgar. Dissolver em ácido. Se o suporte físico fosse destruído, o ciclo poderia parar. Mas e o segundo diário? E se a destruição do diário original libertasse a criatura em vez de aprisioná-la?

\noindent \textbf{Opção 3: Aceitar.}

\noindent Aceitar que o ciclo agora fazia parte dela. Aceitar que o olho direito estava mudando. Aceitar que o cheiro não ia embora. Aceitar que ela era, a partir daquele momento, a nova guardiã do diário --- e que, cedo ou tarde, teria que passá-lo adiante.

\noindent \textbf{--- Mas havia uma quarta opção.}

\noindent Ler o segundo diário. Não como reação ao medo — como ato de investigação. Escolher ativamente entrar no ciclo em vez de ser arrastada por ele. Recusar o papel de vítima e assumir o de testemunha. A diferença, Lúcia pensou, era a diferença entre afogar-se e mergulhar.

\vspace{0.5em}\noindent\rule{\textwidth}{0.5pt}\vspace{0.5em}

\noindent A resposta veio às 3:14 da manhã.

\noindent O cheiro acordou ela.

\noindent Não era o cheiro adocicado de antes. Era mais forte. Mais denso. Cheiro de terra molhada. De osso queimado. De coisa saindo do chão.

\noindent Lúcia abriu os olhos.

\noindent O segundo diário estava aberto sobre a mesa.

\noindent Ela não lembrava de ter aberto.

\noindent A página mostrava uma frase, escrita com a mesma letra trêmula:

\noindent \textit{\textquotedblleft{}Ela já escolheu por você. Você só ainda não sabe.\textquotedblright{}}

\noindent Lúcia fechou o diário. Guardou os dois na mochila. Pagou o hotel. Pegou o ônibus das 5 da manhã de volta para Belo Horizonte.

\noindent Durante toda a viagem, o olho direito ardeu.

\noindent Como se algo estivesse crescendo dentro dele.

\noindent Como se algo estivesse \textit{aprendendo a ver}.

\noindent ---

\noindent\textit{\textbf{↪ Rotas:} \rota{LUC-06} \rotasep \rota{LUC-Entrevista}}

\vspace{0.5em}\noindent\rule{\textwidth}{0.5pt}\vspace{0.5em}

\subsection*{ESCALA DE CONTAMINAÇÃO --- ATUALIZAÇÃO 1}

\noindent \textit{O arquivo detectou que você acompanhou a investigação de Lúcia até o momento da escolha.}

\begin{moltabletwo}
\hline
Pontos & Diagnóstico \\
\hline
0--1 & Leitor intacto. Improvável, mas possível. \\
2--4 & Leitor em exposição. Você já leu documentos que não deveria. \\
5--7 & Leitor contaminado. Lúcia escolheu. Você também já fez a sua escolha --- mesmo que não saiba qual foi. \\
8--10 & Paciente 1.263 em formação. A criatura já tem seu nome. \\
\hline
\end{moltabletwo}

\begin{quote}
Nota do arquivo: os pontos não são cumulativos. A escala se ajusta conforme o arquivo se fecha ao seu redor.
\end{quote}

\vspace{0.5em}\noindent\rule{\textwidth}{0.5pt}\vspace{0.5em}